\section{Expected Research Contributions}

EVERYTHING IS A I HERE, NEEDS REWRITING

The proposed Dual-System Human-Centric Recommender System (DS-HCRS) introduces several conceptual and methodological contributions to the design of recommender systems.
While the implementation and empirical evaluation are planned for future work, the proposed framework outlines a structured approach for incorporating cognitive modeling and long-term optimization into recommendation pipelines.

The expected contributions of this work can be summarized as follows.

\subsection{Dual-System Recommendation Architecture}

This work proposes a dual-system recommender architecture inspired by cognitive theories of human decision-making.
The framework separates recommendation strategies into two complementary components:

\begin{itemize}
    \item \textbf{System 1}, which captures short-term behavioral preferences and predicts immediate user interactions.
    \item \textbf{System 2}, which evaluates recommendations using longer-term utility signals and broader contextual information.
\end{itemize}

By explicitly separating impulsive interaction prediction from reflective optimization, the proposed system introduces a structured way to balance short-term engagement with long-term user outcomes.
This architectural design provides a flexible framework that can integrate different recommendation models while maintaining a clear conceptual separation of objectives.

\subsection{Fatigue-Aware Recommendation Modeling}

Another key contribution of this work is the introduction of an explicit \textit{user fatigue and cognitive state model} within the recommendation pipeline.

Traditional recommender systems assume that additional user interactions always indicate positive engagement.
However, prolonged sessions and repeated exposure to similar content may lead to boredom or cognitive overload.
The proposed fatigue modeling component attempts to estimate this state using observable interaction signals such as repetition, novelty, and session length.

By incorporating fatigue estimates into the recommendation process, the system can adapt its behavior when signs of declining user interest are detected.

\subsection{Adaptive Strategy Selection Through Gating}

The proposed system introduces an adaptive gating mechanism that dynamically determines how recommendations from different components should be combined.

Instead of relying on a single recommendation strategy for all situations, the system adjusts the influence of System 1 and System 2 based on the user's current interaction context.
This adaptive decision process allows the recommender to maintain efficiency during normal usage while introducing more reflective reasoning when the user's experience begins to deteriorate.

Such adaptive control mechanisms provide a potential pathway toward more context-aware recommender systems.

\subsection{Integration of Long-Term Utility Optimization}

The proposed framework incorporates long-term optimization principles into the recommendation pipeline by modeling recommendations as sequential decisions that influence future user behavior.

By combining short-term engagement signals with longer-term indicators such as retention and fatigue, the framework outlines a multi-objective reward structure that can guide recommendation policies toward more sustainable outcomes.

This integration highlights how reinforcement learning-based ideas can be used to move beyond purely engagement-driven recommendation objectives.

\subsection{Human-Centered Evaluation Perspective}

Finally, this work proposes an evaluation framework that extends beyond traditional recommendation accuracy metrics.
In addition to standard ranking measures, the evaluation strategy includes metrics related to diversity, novelty, fatigue indicators, and long-term engagement.

This broader evaluation perspective aligns the system with recent research directions in human-centered recommender systems and responsible AI.
By incorporating metrics related to user experience and cognitive load, the proposed evaluation framework aims to provide a more comprehensive understanding of recommender system behavior.

Overall, the proposed framework contributes a structured approach for combining behavioral prediction, cognitive modeling, and long-term optimization within a single recommender system architecture.
The design emphasizes adaptability and human-centered objectives while remaining compatible with existing recommendation techniques.

